\gddchapter{Research objectives \& Participant information}


\section{Research objectives}

\subsection{Problem statement}

When people are  playing a really well made game some of them  describe feeling like they're directly inside of the digital experience.
Video games vary in their degree in which they can achieve that. This thesis explores weather the voice input modality can have a positive impact on said effect. 
While commands do exist in certain games but they're often require memorisation of the right commands.
These tests attempt to verify if integrating a two way communication between the game engine and a large language model can possibly lead to a deeper immersion compared to traditional input methods.

\subsection{Research questions}

\begin{enumerate}
    \item How would the participants describe the experience of using voice input in the presented
game environment?
    \item How do users feel about using the voice modality compared to traditional input?
\end{enumerate}

\section{Ethical checklist}
    \begin{todolist}
    %   \item[\done] Informed consent
    \item[\done] Informed consent - The participant is consenting to the interview
    \item[\done] Recording consent - The user agrees to the interview being recorded
    \item[\done] Anonymity choice - real name or a pseudonym (reflected on the main page)
    \item[\done] Data storage - Confirm the data storage policy with the participant
    \end{todolist}

\section{Participant information}
\begin{table}[h]
\centering
\begin{tabular}{|p{5cm}|p{10cm}|}
\hline
\textbf{Field} & \textbf{Response} \\ \hline
Participant ID / Name & Szymon Wasyl \\[1.5cm] \hline
Gender & Male \\[1.5cm] \hline
Age & 26 \\[1.5cm] \hline
Primary language & Polish \\[1.5cm] \hline
English proficiency & B2/C1 \\[1.5cm] \hline
Native language & Polish \\[1.5cm] \hline
Interview language & English \\[1.5cm] \hline
Gaming Experience (e.g., Frequent, Casual, Non-gamer) & Retired serious gamer, rpg player these days \\[1.5cm] \hline % Fixed here
Familiarity with Voice Assistants (e.g., Alexa, Siri) & No usage  \\[1.5cm] \hline
Physical Accessibility Context & Fully able-bodied \\[2cm] \hline
\end{tabular}
\end{table}


\section{Researcher subjectivity statement}

\subsection{Personal interest}
%  How has your background as a developer or your interest in handicapped accessibility led you to this topic?

    As a person which suffer from impaired vision I'm more inclined to explore the topics that are centered around accessibility than a regular able bodied person.
Another angle which influences my personal affection for this topic is the fact that I've recieved 3 year eductaion in Video Game design which put a lot of emphasis on immersion
as a valuble part of game creation. These two factors naturally led me to explore this topic.
\subsection{Preconcivied notions}
% What are your current "presuppositions about knowledge" in this field? For example, do you already believe voice control will be "more personal" but "slower" than a keyboard?
My presumptions are that the voice control might be feel more immersive than using traditional voice input. This originates from my experience in playing pen and paper roleplaying 
games in which all of the navigation and controls are done by voice. In my personal opinion the levels of immersion in such experience is unparalelled. This immersion stems not only from the voice 
modality since there's also the imagination factor (players imagine the game instead of seeing the screen). It's hard to say how much the voice input influences the immersion, but my hunch is that there is a connection between them.

Another assumption that I have is that the navigation will be slower than the keyboard and potentially more frustrating. Keyboard input is predictable and more linear (since typically there is only one way to issue any given command in a game by using a keyboard), where as 
voice control depends heavily on the quality of integration into game engine, the voice processing technology, the power of the backend processing, 
the quality of microphone, room acoustics and even the accent of the speaker. All of those factors make me believe that there's a much higher potential for slower and more irritating
navigation compared to keyboard based input.  
\subsection{Interviewer stance}
% Acknowledge your role as the creator of the prototype. How might your desire for the WhisperCPP backend to succeed influence how you probe for "frustration" or "success"?
I understand that I'm inherently biased as a creator of the prototype. Every creator wishes for their creation to succeed. This bias in itself is inherently dangerous, since it can easily taint the testing process, arriving at conclusion and making the testing data fit that preconcieved narrative. In a best case scenario my work should be tested by a fully independent researcher, but I can still do my best to protect myself from this bias. I plan on remining myself as much as possible about this bias during testing and I will conduct self reflection before and after testing searching trying to establish if I had fallen for it during the session (and if so I;ll note those moments and make sure to not repeat them). 
